% =====================================================================
\section{SPI: быстрая периферия и SD-карты}
% =====================================================================
\lessonintro
  {подключать быстрые устройства по SPI и безопасно делить одну шину между несколькими микросхемами}
  {урок 7 (I\textsuperscript{2}C), урок 2 (провода и логические уровни)}
  {плата сохраняет журнал событий на SD-карту}
  {модуль microSD с логикой 3,3 В, карта памяти, короткие провода}

\subsection{Чем SPI отличается от I\textsuperscript{2}C}

SPI --- синхронная шина без адресов. Ведущий (ESP32-S3) создаёт тактирование \code{SCK}; линии
\code{MOSI} и \code{MISO} несут данные к устройству и от него. Для каждого ведомого нужен свой
\textbf{CS} (Chip Select): только устройство с активным CS имеет право отвечать.

\begin{center}\small
\begin{tabularx}{\textwidth}{@{}l X X@{}}
\toprule
 & I\textsuperscript{2}C & SPI \\
\midrule
Провода & SDA, SCL; адрес в пакете & SCK, MOSI, MISO + отдельный CS на устройство \\
Скорость & обычно 100 кГц--1 МГц & единицы--десятки МГц \\
Сильная сторона & много простых датчиков & дисплеи, SD-карты, АЦП с быстрым потоком данных \\
\bottomrule
\end{tabularx}
\end{center}

\begin{caution}
SPI не задаёт единой распиновки. В Arduino-ESP32 сигналы можно назначить программно, но у модуля
SD и платы должны совпадать \emph{все четыре} названия. Питание и логика также только 3,3 В.
\end{caution}

\subsection{Одна шина, несколько устройств}

К линиям SCK, MOSI и MISO параллельно подключают все устройства. CS у каждого свой. Перед обменом
библиотека опускает нужный CS в LOW, после --- возвращает HIGH. Если два CS активны одновременно,
два выхода MISO могут физически спорить друг с другом: результатом будут ошибки или повреждение.

\begin{figure}[H]
\centering
\begin{tikzpicture}[font=\sffamily\scriptsize]
  \node[blk,fill=blkcpu,minimum width=2.4cm,minimum height=3.6cm,anchor=north west] (m) at (0,1.6) {ESP32-S3\\ ведущий};
  \foreach \y/\n/\c in {1.2/SCK/esp,0.7/MOSI/accent,0.2/MISO/okgreen}
    \draw[thick,\c] (2.4,\y) -- node[pos=0.07,above,text=\c]{\n} (10.4,\y);
  \foreach \cx/\name/\cs/\mx/\yy in {5.2/microSD/CS1/1.6/-2.55, 8.8/TFT-дисплей/CS2/2.0/-2.95}{
    \node[blk,fill=blkper,minimum width=2.2cm,minimum height=1.1cm] at (\cx,-1.6) {\name};
    \foreach \dx/\y/\c in {-0.6/1.2/esp,-0.2/0.7/accent,0.2/0.2/okgreen}{
      \draw[thick,\c] (\cx+\dx,\y) -- (\cx+\dx,-1.05);
      \fill[\c] (\cx+\dx,\y) circle (1.6pt);}
    \draw[thick,black!70] (\cx+0.6,-1.05) -- (\cx+0.6,-0.5) node[above,text=black]{CS};
    \draw[thick,black!70] (\cx,-2.15) -- (\cx,\yy) -- (\mx,\yy) node[pos=0.5,below]{\cs} -- (\mx,-2.0);
  }
\end{tikzpicture}
\caption{Общая шина SPI: линии SCK, MOSI и MISO идут ко всем устройствам, а линия выбора CS у каждого
устройства своя. В каждый момент активен (LOW) только один CS}
\end{figure}

Выберите, например, GPIO12 (SCK), GPIO11 (MOSI), GPIO13 (MISO), GPIO10 (CS SD), если они свободны
на вашей ревизии. Перед подключением всегда сверяйтесь с pinout именно своего модуля: у плат с PSRAM
часть выводов занята памятью.

\begin{figure}[H]
\centering
\begin{circuitikz}[scale=0.95,transform shape]
  \node[blk,fill=blkcpu,minimum width=2.8cm,minimum height=6.2cm] (m) at (0,0) {};
  \node[blk,fill=blkper,minimum width=2.8cm,minimum height=6.2cm] (d) at (8,0) {};
  \node[font=\sffamily\small\bfseries] at (0,3.7) {ESP32-S3};
  \node[font=\sffamily\small\bfseries,align=center] at (8,3.85) {Модуль microSD\\[-2pt]\scriptsize логика \SI{3.3}{V}};
  \foreach \y/\a/\b/\c in {2.5/3V3/VCC/wred, 1.5/{\pin{10}}/CS/wblue, 0.5/{\pin{11}}/MOSI/wgreen,
      -0.5/{\pin{12}}/SCK/wyellow, -1.5/{\pin{13}}/MISO/worange, -2.5/GND/GND/wblack}{
    \node[font=\sffamily\scriptsize,anchor=east] at (1.4,\y) {\a};
    \node[font=\sffamily\scriptsize,anchor=west] at (6.6,\y) {\b};
    \draw[thick,\c] (1.4,\y) -- (6.6,\y);
  }
\end{circuitikz}
\caption{Подключение модуля microSD к SPI-шине: \pin{12} --- SCK, \pin{11} --- MOSI, \pin{13} --- MISO,
\pin{10} --- CS. Питание \SI{3.3}{V} и общая земля. Порядок контактов на вашем модуле может отличаться:
ориентируйтесь на надписи}
\end{figure}

\subsection{Первый журнал на SD-карте}

\codecap{Инициализация SPI и запись строки в файл}
\codeinput{l10_spi_sd}

Начинайте с \SI{4}{MHz}. Макетная плата, длинные провода и плохое питание добавляют паразитную
ёмкость и искажают фронты. Лишь после стабильной работы поднимайте частоту небольшими шагами.
Частота \SIrange{20}{40}{MHz}, которую часто обещает модуль, обычно достижима только на короткой
разводке платы.

\subsection{Разбор ошибки: «карта то работает, то нет»}

Проверьте в таком порядке:
\begin{enumerate}
  \item Общая земля, питание 3,3 В и настоящий преобразователь уровней на модуле SD. Модули,
        рассчитанные на Arduino Uno, иногда требуют 5 В питания, но их входы всё равно нельзя
        считать безопасными для ESP32 без проверки схемы.
  \item CS: он не должен быть оставлен «в воздухе»; при старте задайте его OUTPUT и HIGH.
  \item Распиновка MOSI/MISO: они легко меняются местами. MISO идёт \emph{к} ESP32.
  \item Понизьте SPI-частоту до \SI{1}{MHz}, укоротите соединения и добавьте развязывающий
        конденсатор \SI{0.1}{\micro F} около модуля.
  \item Закрывайте файл после записи. Иначе сообщение остаётся в буфере и пропадёт при отключении питания.
\end{enumerate}

\begin{summary}
\begin{itemize}[leftmargin=*]
  \item SPI быстр, но требует отдельного CS для каждого ведомого устройства.
  \item На макетной плате сначала выбираем низкую частоту и короткие провода.
  \item Ошибка SD чаще всего означает питание, CS, перепутанные линии или незакрытый файл.
\end{itemize}
\end{summary}

\begin{selfcheck}
\begin{enumerate}
  \item Зачем каждому SPI-устройству своя линия CS?
  \item Почему на макетной плате начинают с низкой частоты SPI?
  \item Почему файл обязательно закрывают после записи?
\end{enumerate}
\end{selfcheck}

\begin{tasks}
\begin{enumerate}
\item \lv{1} Нарисуйте таблицу выводов SPI вашей платы: MOSI, MISO, SCK, CS.
\item \lv{1} Запустите пример с SD-картой и выведите её размер и тип.
\item \lv{2} Записывайте в \code{/log.csv} время работы и показания потенциометра, добавляя заголовок только в новый файл.
\item \lv{2} Прочитайте файл обратно и выведите его в Serial Monitor построчно.
\item \lv{3} Подключите TFT-дисплей и SD-карту к общей шине: нарисуйте таблицу CS и проверьте, что они не мешают друг другу.
\item \lv{4} Найдите максимальную надёжную частоту SPI для своей сборки, увеличивая её и записывая результаты в таблицу.
\end{enumerate}
\end{tasks}
